\section{Relative Stokes with common tangential strains}
\label{sec:r13-common-strain-stokes}

The local equation relevant to a positive face involves primitives on
all incident half-products.  Equal tangential strains suffice for its
boundary calculation.  Their tangential primitive traces may differ
by affine Killing fields.  Keeping these fields in the Green formula
gives the full relation Hessian at a completed curve star.

\subsection{The half-products and the boundary tensor}

Let $T$ be a real vector space of dimension $d\geq1$, with fixed
lattice measure, and let $B\subset T$ be a bounded convex open
polytope.  Let $N$ be a finite-dimensional real vector space and
let
\[
 \Sigma=\sum_{i\in I}w_i[\mathbb R_{\geq0}v_i],\qquad
 v_i\ne0,\qquad \sum_iw_iv_i=0,
 \tag{CS1}
\]
be a finite labelled balanced one-cycle.  Each $v_i$ is primitive
for its specified ray lattice.  The real weights may be signed in
the real statement.  Distinct labels are retained even when their
vectors coincide.  Put $W_i=T\oplus\mathbb Rv_i$.  On the closed
half-product $B\times[0,\epsilon)v_i$ let
\[
 b_i(t,s)=\sum_{a=1}^d b_{i,a}(t,s)\,dt_a+c_i(t,s)\,ds
 \tag{CS2}
\]
be continuous and finite piecewise smooth, with smooth extensions
near the closed cells and locally bounded first derivatives.  All
coefficient walls are included, including walls meeting $s=0$.
Write $A_i=e(b_i)=\operatorname{sym}(db_i)$, using distributional
derivatives.  Continuity makes its first derivative the bounded
cellwise derivative, with no coefficient-jump measure.

Assume the tangential boundary covectors have one common strain:
\[
       e_T\bigl(b_i(t,0)|_T\bigr)=A_T
       \quad\hbox{as distributions on }B\quad(i\in I).
 \tag{CS3}
\]
Their values need not coincide.  Define the normal flux and the
$N$-valued tangential covector
\[
 z(t)=\sum_iw_ic_i(t,0),\qquad
 B_a(t)=\sum_iw_ib_{i,a}(t,0)v_i.
 \tag{CS4}
\]

\begin{lemma}[The unmatched trace is a Killing tensor]
\label{lem:r13-weighted-killing-trace}
For every $\lambda\in N^*$, the covector
$\sum_a\lambda(B_a)dt_a$ is an affine Killing covector on $B$.
Equivalently,
\[
 \partial_aB_b+\partial_bB_a=0
       \quad\hbox{distributionally for all }a,b.
 \tag{CS5}
\]
For any fixed index $i_0$, each difference
$K_i=b_i(\cdot,0)|_T-b_{i_0}(\cdot,0)|_T$ is likewise affine
Killing.
\end{lemma}

\begin{proof}
Apply $e_T$ to CS4.  By CS3 the result is
$A_T\otimes\sum_iw_iv_i=0$.  This proves CS5.  For completeness,
if $k_a$ are distributional coefficients satisfying
$\partial_ak_b+\partial_bk_a=0$, then
\[
 2\partial_a\partial_bk_c
 =\partial_a(\partial_bk_c+\partial_ck_b)
  +\partial_b(\partial_ak_c+\partial_ck_a)
  -\partial_c(\partial_ak_b+\partial_bk_a)=0.
 \tag{CS6}
\]
Thus every coefficient is affine on the connected convex set
$B$, and its linear part is alternating.  This conclusion applies
to $\lambda(B)$ and, by CS3, to every $K_i$.  In particular it
already includes all tangential coefficient walls.
\end{proof}

\begin{theorem}[Common-strain product Stokes]
\label{thm:r13-common-strain-stokes}
For every smooth compactly supported ambient function $\chi$ on
a neighborhood of the truncated product and every constant
symmetric tangent form $\Omega$ of bidegree $(d-1,d-1)$,
\[
 \sum_iw_i\int_{B\times[0,\epsilon)v_i}
                       A_i\wedge D\chi\wedge\Omega
       =-\int_B z\,D_T\chi_0\wedge\Omega,
 \qquad \chi_0(t)=\chi(t,0).
 \tag{CS7}
\]
The support of $\chi$ stays away from the outer boundaries.
Here $D=d'd''$, with the positive product orientation and the
original primitive ray parameters.
\end{theorem}

\begin{proof}
Let the symmetric matrix $C^{ab}$ represent the tangent dual
$\Omega$.  The complete half-product Green formula FS6 gives,
on the branch $i$,
\[
 \int A_i\wedge D\chi\wedge\Omega
 =\int_B\sum_{a,b}C^{ab}
 \bigl(b_{i,a}(t,0)\,\partial_b\partial_s\chi_i(t,0)
              -c_i(t,0)\,\partial_a\partial_b\chi_0(t)\bigr)dt,
 \tag{CS8}
\]
where $\chi_i(t,s)=\chi(t,sv_i)$.  Full continuity of $b_i$
on its half-product removes every internal-wall term in that
formula.  Its integrand before taking the boundary is the full
contraction
\[
 \sum_{a,b}C^{ab}
   \bigl((A_i)_{ab}(\chi_i)_{ss}
          +(A_i)_{ss}(\chi_i)_{ab}
                     -2(A_i)_{as}(\chi_i)_{bs}\bigr).
 \tag{CS9}
\]

Put $\xi(t)=d_N\chi(t,0)\in N^*$.  The sum of the first
boundary terms in CS8 is
\[
       \int_B\sum_{a,b}C^{ab}
                       \langle\partial_b\xi(t),B_a(t)\rangle dt.
 \tag{CS10}
\]
Integrate once in $t$.  This is legitimate distributionally,
so no tangential wall term is omitted.  Since $C$ is symmetric,
CS5 makes the result
\[
 -\frac12\int_B\sum_{a,b}C^{ab}
                   \langle\xi,\partial_bB_a+\partial_aB_b\rangle dt
 =0.
 \tag{CS11}
\]
The remaining terms in CS8 are precisely the right side of
CS7.  Thus the tangent cancellation occurs after integration;
it requires the common tangential strain in CS3, rather than
a common tangential primitive value.
\end{proof}

\subsection{Intrinsic Hessians and the complete relation space}

The normal flux depends on a quotient section, whereas its full
tangential Hessian does not.  To state this, start with an ambient
space $W$ and its tangent subspace $T$.  The normal vectors are
in $W/T$, and a common linear section identifies them with the
$v_i$ used above.  The incident space is
$W_i=T+\mathbb Rv_i$.

\begin{lemma}[Section and primitive-gauge invariance]
\label{lem:r13-flux-hessian-invariance}
Under CS3, changing the common linear section of $W\to W/T$
changes $z$ by an affine function on $B$.  Changing any branch
primitive by a strain-preserving affine Killing covector also
changes $z$ by an affine function.  Consequently
$\operatorname{Hess}_Tz$ is independent of these choices.
\end{lemma}

\begin{proof}
Two sections differ by a linear map $L:W/T\to T$.  The new
normal lift is $v_i+Lv_i$, so
\[
 z'-z=\sum_iw_i b_i(t,0)|_T(Lv_i)
      =\sum_iw_iK_i(t)(Lv_i).
 \tag{CS12}
\]
In the second equality the common $b_{i_0}$ contribution vanishes
by balancing.  Each $K_i$ is affine by the preceding lemma;
hence CS12 is affine.  The same proof allows individual changes
$v_i\mapsto v_i+t_i$ whose weighted sum $\sum_iw_it_i$ is zero.
For a strain-preserving change on a branch, CS6 in the full
half-product variables gives an affine Killing field.  Its
normal evaluation along $s=0$ is affine in $t$, proving the
second assertion.
\end{proof}

Define the complete labelled normal relation space and raw
incident relation space by
\[
 R_\Sigma=\ker\left(\mathbb R^I\longrightarrow N,
                         (a_i)\longmapsto\sum_ia_iv_i\right),
 \quad
 \mathcal K=\ker\left(\bigoplus_iW_i\longrightarrow T\oplus N\right).
 \tag{CS13}
\]
For $a\in R_\Sigma$ and $u=(u_i)\in\mathcal K$ put
\[
 z_a(t)=\sum_i a_i b_i(t,0)(v_i),\qquad
 \mathcal F_u(t)=\sum_i b_i(t,0)(u_i).
 \tag{CS14}
\]
Every sum map here retains simultaneous relations among all
incident spaces.  It is not replaced by pairwise intersections.

\begin{proposition}[The exact Hessian reduction]
\label{prop:r13-raw-hessian-reduction}
Under CS3, the full Hessian of every raw defect vanishes if and
only if $\operatorname{Hess}_Tz_a=0$ for every
$a\in R_\Sigma$.  More precisely, write uniquely
$u_i=t_i+a_iv_i$.  Then $\sum_it_i=0$, $a\in R_\Sigma$, and
\[
        \mathcal F_u(t)=z_a(t)+\sum_iK_i(t)(t_i).
 \tag{CS15}
\]
The second term is affine.  The correspondence therefore
identifies the effective full relation Hessian with the linear
map
\[
 R_\Sigma\longrightarrow\mathcal D'(B,\operatorname{Sym}^2T^*),
                    \quad a\longmapsto\operatorname{Hess}_Tz_a.
 \tag{CS16}
\]
\end{proposition}

\begin{proof}
Use the fixed splitting $T\oplus N$.  The condition
$\sum_iu_i=0$ separates into the two assertions in the
statement.  Subtracting $b_{i_0}(\sum_it_i)=0$ yields CS15.
Conversely every $a\in R_\Sigma$ is represented by the raw
relation $u_i=a_iv_i$.  This proves both directions and the
linear description CS16.  In particular, tangential duplicate
relations contribute only affine functions even when their
primitive traces have not been aligned.
\end{proof}

\begin{corollary}[One balanced relation gives the full Hessian]
\label{cor:r13-single-balance-full-hessian}
Assume $R_\Sigma=\mathbb Rw$, and suppose
\[
 \sum_iw_i\int A_i\wedge D\chi\wedge\Omega=0
       \quad\hbox{for all }\chi,\Omega\hbox{ in CS7}.
 \tag{CS17}
\]
Then every $\mathcal F_u$ has zero full distributional Hessian
on $B$.  Equivalently, every raw relation defect is affine on
$B$, including across its coefficient walls.
\end{corollary}

\begin{proof}
Choose $\chi(t,n)=\varphi(t)\psi(n)$, where $\varphi$ is
arbitrary compact smooth on $B$ and $\psi=1$ near zero.
CS7 and CS17 give
$\int_Bz\,D_T\varphi\wedge\Omega=0$.  All constant symmetric
tangent duals occur, so this is
$\operatorname{Hess}_Tz_w=0$.  Since $w$ spans the entire
$R_\Sigma$, the preceding proposition proves the assertion.
The joint test symbols are the actual pairs
$w\otimes C$, for every symmetric $C$; they span
$R_\Sigma\otimes\operatorname{Sym}^2T$.
\end{proof}

More generally, if actual complete test constructions supply
CS17 for balanced vectors $w^{(1)},\ldots,w^{(r)}$ on the
same coefficient system, their conclusions detect exactly the
restriction of CS16 to their linear span.  Equality of that
span with $R_\Sigma$ is the required finite test.  A list of
separately spanning scalar functions does not supply this
condition on the actual balanced vectors.

\subsection{Normalized marked curves}

Let $C/k$ be a proper smooth integral curve over the original
finite residue field.  Retain its full constant field and put
$\delta_p=[\kappa(p):k]$.  For a finite marked set $M$ with
$|M|\geq2$,
let $g_1,\ldots,g_m\in k(C)^\times$ have divisor supports
in $M$, and let
\[
 \nu_p=(\operatorname{ord}_p g_j)_j\ne0,
 \qquad \nu_p=m_pv_p,
 \qquad m_p\in\mathbb Z_{>0},\quad v_p\text{ primitive}.
 \tag{CS18}
\]
Assume the marked divisor rows of the chosen functions span
the degree-zero hyperplane
\[
 \operatorname{span}_{\mathbb R}
       \{(\operatorname{ord}_pg_j)_{p\in M}:1\leq j\leq m\}
             =\delta^\perp\subset\mathbb R^M.
 \tag{CS19}
\]
The retained finite-field Picard construction supplies such
functions.  If several prescribed curves lie on one supplied
smooth projective integral model, ML supplies their exact
simultaneous restrictions in the same ambient function field.

\begin{lemma}[Primitive conversion and complete relation kernel]
\label{lem:r13-curve-primitive-kernel}
For $w_p=\delta_pm_p$, the labelled column map
$(a_p)\mapsto\sum_pa_pv_p$ has kernel $\mathbb Rw$.
For the normalized curve parameter $u=-\log|z_p|$ and primitive
image parameter $s=m_pu$, a logarithm with curve order
$\operatorname{ord}_p g$ has primitive slope
\[
 h_p(g)=\frac{\operatorname{ord}_p g}{m_p},
      \qquad w_ph_p(g)=\delta_p\operatorname{ord}_p g.
 \tag{CS20}
\]
\end{lemma}

\begin{proof}
The matrix with columns $\nu_p$ has kernel $\mathbb R\delta$
by CS19.  Since its columns are $m_pv_p$, rescaling the columns
gives the asserted kernel.  Substituting $u=s/m_p$ in the
linear branch logarithm proves CS20.  In particular all marks
remain labelled; CS19 also excludes a zero column.  Distinct
marks cannot yield the same primitive ray, since their
coordinate difference would be a kernel vector with entries
of both signs, whereas all $w_p$ are positive.
\end{proof}

The equality $w_p=\delta_pm_p$ in an analytic application is
the actual source calibration, to be supplied by its ordinary
comparison theorem.  CS20 describes its conversion and does
not replace the calibration.  The tangent-face lattice is
the saturated actual lattice.  Any index of its projection
to an older, nonsaturated tangent face remains in that separate
projection formula.

\subsection{Original ordinary tests and tangent-only sources}

We state the analytic application with a finite actual source
presentation.  Let $U$ be an ordinary original analytic
neighborhood, and suppose a fixed finite moment refinement has
a protected calibrated product presentation over $B$ with
labelled main cycle $\Sigma$ as in CS1.  Every added complete
source cycle is balanced, and its old coefficient image is
contained in the same tangent plane.  The retained tangent
coordinates split the actual tangent lattice.  The represented
tensor on such an added source is the literal pullback of
$A_T$.  Require that all compact ambient tests
$\chi\Omega$ used above are realized as original compact
smooth CLD tests on $U$, and that integration uses the complete
source sum with its actual weights.  These are the ordinary
protection and source-comparison hypotheses.

\begin{theorem}[Actual raw-Hessian detector on a completed curve block]
\label{thm:r13-actual-completed-curve-detector}
Suppose the main branch primitives satisfy CS2--CS3, and their
strains represent an ordinary current $Z$ for which
$D[Z]=0$ on every original compact complementary smooth CLD
test on $U$.  If its calibrated main curve star satisfies
CS18--CS19 with actual weights $w_p=\delta_pm_p$, then every
raw incident defect has zero full distributional Hessian on
$B$.  No tangential primitive alignment is needed before this
calculation.
\end{theorem}

\begin{proof}
Fix one main tangent trace $b_0$ with $e_Tb_0=A_T$.  On every
additional source use its pullback.  Because the old coefficient
image lies literally in the tangent plane, its strain is the
represented pullback of $A_T$.  In the quotient splitting given
by the retained tangent projection, its normal trace is zero.
The common $b_0$ tangent boundary terms cancel on that whole
balanced source cycle.  If a different quotient section is
used, the lemma on section changes gives the same Hessian
statement.  All source labels and their absolute carriers have
remained in the integral.

The original current equation and the complete source formula
therefore give CS17 for the main cycle.  Its primitives only
need CS3, by the common-strain Stokes theorem.  The curve
primitive-kernel lemma gives $R_\Sigma=\mathbb Rw$, so the
single-balance corollary yields the full raw relation Hessian.
Each symmetric tangent dual is tested together with this same
complete actual one-cycle.  All internal coefficient walls are
already included in CS8--CS11.
\end{proof}

On a positive codimension-one face of a three-dimensional normal
fan, the base $B$ includes both old tangent variables and the
two directions along that face.  The theorem thus tests the
entire Hessian in those variables, including angular and mixed
entries.  Suppose these local Hessian identities cover the whole
connected convex product of the old tangent base with the
relative interior of the positive face, using one fixed finite
coefficient system.  If that system extends continuously to the
central normalized face, the affine expressions agree on
overlaps.  The zero-central-value condition then removes their
constant and old tangent-linear terms.  This gives JF7 on that
whole face.  A single compact box away from the central face
only gives the local affine conclusion of the theorem.  The finite
skew-gauge membership JF8 and the raw systems on other faces
remain distinct finite equations.

\paragraph{Proof inputs.}
The half-product identity is the reviewed FS6--FS9 in
\texttt{finite\_ps\_relative\_stokes.tex}, retained in Round 12
under the frozen Round 11 inputs.  The raw relation convention
is JF3 of Round 12.  The finite Picard row construction is
PP7--PP11 of the retained pointed Picard theorem; the finite
simultaneous ambient restriction theorem is ML of Round 12.
The new deductions here are the cancellation CS10--CS11 with
unequal tangential primitive traces, the section-invariant full
relation Hessian CS12--CS16, and their actual completed-curve
combination under the stated complete source hypotheses.
